\section{Discussion}

\subsection{The theorem in one sentence}

The finite result of Paper 14 can now be stated as follows:

\begin{quote}
Admissible native transport reduces lawful future freedom from
\(8\) to \(2\) to \(1\), while contextual distinction remains
retained; the resulting designation response selects exactly the
ten-dimensional five-address-hidden sector, and the Thalean
quadratic independently contracts centered registered contrast
without finite-stage erasure.
\end{quote}

The important word is \emph{independently}.

Native relational loading and quadratic contrast concentration are
not two descriptions of the same operation.

They are complementary finite mechanisms.

\subsection{Selection, retention, and concentration}

The mature architecture separates three operations.

\paragraph{Selection.}
For a live prefix \(p\), let
\[
\mathcal C(p)
\]
be the lawful completion family.

Native loading is ordered by
\[
p\preceq p'
\quad\Longrightarrow\quad
\mathcal C(p')\subseteq\mathcal C(p).
\]

For the certified carrier,
\[
8\longrightarrow2\longrightarrow1.
\]

Thus admitted transport progressively selects a future.

\paragraph{Retention.}
Future uniqueness does not collapse the complete relational state.

At the two-edge threshold there remain \(120\) distinct contexts
with \(120\) distinct completed words.

The native designation response
\[
\Delta
\]
acts entirely in the retained five-address-hidden sector:
\[
\Delta^{\mathsf T}\Delta
=
\Delta\Delta^{\mathsf T}
=
\frac34H,
\]
where
\[
\operatorname{rank}(H)=10.
\]

The declared five-address readout satisfies
\[
\rho\Delta=0.
\]

Thus the retained response is nontrivial internally while invisible
to the coarse readout.

\paragraph{Concentration.}
The normalized quadratic
\[
\widehat Q=\frac1{98}Q
\]
fixes common mode and contracts every centered eigendirection.

Its sharp centered factor is
\[
\frac9{49}.
\]

Consequently,
\[
R(\widehat Q^{\,k}x)
\le
\left(\frac9{49}\right)^k R(x),
\]
while every finite iterate remains invertible.

Selection, retention, and concentration are therefore distinct
finite operations with exact algebraic relations.

\subsection{Native loading sharpens rather than smooths}

The canonical completion-occupancy register provides a decisive
test of the relation between native loading and Paper 14 contrast.

Its exact contrast values are
\[
\zeta:
\frac{31}{55}
\longrightarrow
\frac{41}{65}
\longrightarrow
\frac{7}{10}.
\]

Likewise,
\[
\rho^2:
\frac{31}{24}
\longrightarrow
\frac{41}{24}
\longrightarrow
\frac73.
\]

The common component remains fixed:
\[
C^2=\frac{12}{5}.
\]

All \(60\) strict \(8\to2\) extensions increase both contrast
statistics.

All \(120\) strict \(2\to1\) extensions do the same.

There are no exceptions.

Thus the native loading process is contrast-antimonotone relative to
the normalized quadratic.

The correct interpretation is:

\[
\text{native loading}
\longrightarrow
\text{future sharpening},
\]
whereas
\[
\widehat Q
\longrightarrow
\text{common-mode concentration}.
\]

This distinction prevents a false identification that earlier
versions of the saturation conjecture left open.

\subsection{The retained-sector identity}

The strongest native-to-register crosswalk is not
\[
T\mapsto\widehat Q.
\]

It is
\[
\Delta^{\mathsf T}\Delta=\frac34H.
\]

This identity says that the quadratic response of context-activated
native designation is exactly the projector onto the retained
ten-dimensional relational sector, up to the scalar \(3/4\).

Because also
\[
\Delta\Delta^{\mathsf T}=\frac34H,
\]
the normalized operator
\[
\sqrt{\frac43}\,\Delta
\]
is an isometry on the retained sector.

Thus native contextual transport carries an exact internally
nontrivial response even though the five-address readout sees
\[
\rho\Delta=0.
\]

This is a finite realization of retained distinction behind
registration.

\subsection{Relation to the Thalean response algebra}

The native loading Gram belongs to the same exact centered
transport-response algebra as the centered Thalean quadratic.

If \(R_0,R_1,R_2\) are the three certified centered response classes,
then
\[
\Delta^{\mathsf T}\Delta
=
\frac1{48}R_0
-
\frac1{16}R_1
+
\frac1{32}R_2.
\]

The centered Thalean quadratic satisfies
\[
PQP
=
\frac{139}{90}R_0
-
\frac{28}{15}R_1
+
\frac{19}{60}R_2.
\]

The operators are not proportional.

Their relationship is therefore algebraic rather than identificatory:
native loading and quadratic concentration inhabit the same certified
response algebra while selecting different directions within it.

\subsection{Two distinct saturation notions}

The upgraded theory contains two theorem-level saturation notions.

\paragraph{Registered-future saturation.}
A native prefix is saturated when
\[
|\mathcal C(p)|=1.
\]

This means lawful future freedom has become unique.

It does not mean complete relational distinction has disappeared.

\paragraph{Contrast saturation.}
A registered sequence is contrast-saturating when its declared
contrast approaches a stable minimum while complete state remains
nontrivial.

Normalized quadratic aggregation supplies the exact example
\[
\zeta(\widehat Q^{\,k}x)\longrightarrow0
\]
whenever the input has a nonzero common component.

Thus
\[
\text{registered-future saturation}
\neq
\text{contrast saturation}.
\]

The first is selection of a unique lawful future.

The second is suppression of centered registered contrast.

\subsection{No transfer into the common mode}

Contrast concentration under \(\widehat Q\) must still not be read as
transfer from the centered sector into common mode.

The common and centered subspaces are invariant.

If
\[
P_0x=0,
\]
then
\[
P_0\widehat Q^{\,k}x=0
\]
for every \(k\).

A purely centered input remains centered and contracts toward zero.

The concentration theorem is therefore a statement about differential
spectral persistence.

This remains distinct from the native carrier mechanism, where
transport creates contextual structure and designation becomes
operational after the first admitted step.

\subsection{Invertible stages and projection loss}

Every finite power of \(\widehat Q\) remains invertible.

Yet
\[
\widehat Q^{\,k}\longrightarrow P_0,
\]
and
\[
\operatorname{rank}(P_0)=1.
\]

Thus exact finite-stage distinction can remain recoverable while a
declared contrast readout becomes arbitrarily insensitive to it.

The five-address theorem supplies a stronger exact form of the same
general principle.

The hidden projector \(H\) has rank ten.

States differing entirely within that sector can remain distinct in
the fifteen-state register while having identical five-address
readout.

Projection loss and finite-stage invertibility are therefore two
different mechanisms by which visible individuation can weaken
without complete-state merger.

\subsection{Global sigma is a separate frontier}

The unresolved inter-face sigma remains mathematically interesting,
but it is no longer a missing premise of the finite saturation
theorem.

The native registered carrier used in Sections 8 and 14 is already
an admitted finite transport system with exact successor and
completion laws.

The sigma problem instead asks whether a stronger inter-face seam
transport can be made globally compatible under its declared
four-step closure conditions.

Therefore
\[
\text{sigma open}
\]
does not imply
\[
\text{native saturation open}.
\]

The current sigma frontier remains:

\[
2816
\]
neighbor-frame-\(E\)-closed local options,

\[
376
\]
completed solver iterations,

\[
175069
\]
unique nogoods,

and a best observed partial closure of
\[
32
\]
flags, without either a complete global \(r_3\) or a proof of
infeasibility.

That is a separate global-compatibility problem.

\subsection{Long range as lawful accumulation}

The finite theory still supports the architectural principle that
long-range registered structure should arise from lawful finite
accumulation rather than from an independently inserted long-range
selector.

The exact examples remain:

\[
x\longmapsto Qx
\]
on the sector register,

\[
d\delta=\omega
\]
on the binary closure complex,

\[
d\eta=\Omega
\]
on its integer refinement,

and
\[
\tau_{\mathrm T}(\gamma)
=
\operatorname{wind}_{r_1}(\gamma)
\]
on registered-history cycles.

Their domains and coefficient systems differ.

No theorem identifies them as one universal sum.

\subsection{Clean release and overcompression remain conditional}

The new native loading theorem changes the status of the clean-release
picture.

It is no longer correct to assume that increasing native loading
itself drives centered contrast downward.

For the canonical native completion register, it does the opposite.

A larger cycle containing both selection and concentration might still
possess an intermediate clean-release or overcompression regime, but
such a composition requires an explicit law specifying when and how
the native selection dynamics couple to a contrast-concentrating
operation.

No such cycle is proved in this paper.

Clean release and overcompression therefore remain internal
conditional concepts.

\subsection{Finite capacity remains a separate hypothesis}

The native loading theorem supplies a finite completion family and an
exact finite distinction count.

It does not prove a universal bound on all possible real-valued
relational states.

The fifteen-state sector space remains a real vector space.

Therefore finite graph dimension must not be silently identified with
finite physical information capacity.

Any stronger capacity statement requires a declared notion of
independent distinction, such as finite cardinality, linear rank,
finite-resolution distinguishability, or another explicitly defined
independence structure.

\subsection{No singularity is required internally}

Nothing in the finite theorem requires a divergent variable.

Native future freedom reaches
\[
|\mathcal C|=1
\]
at a finite stage.

The retained response remains finite:
\[
\Delta^{\mathsf T}\Delta=\frac34H.
\]

The normalized quadratic remains finite at every iterate and converges
to the finite projector
\[
P_0.
\]

Thus the internal mathematics reaches saturation through finite
selection, retained distinction, and projection-relative concentration
rather than through divergence.

This is not a theorem about singularities in general relativity.

\subsection{Methodological posture}

The upgraded result preserves the methodological rule:

\begin{quote}
Do not protect the idea. Make the idea protect itself.
\end{quote}

The canonical native loading test could have confirmed the earlier
expectation that loading itself decreases contrast.

It did not.

Every one of the \(180\) strict extensions moved contrast in the
opposite direction.

The correct response is not to change the readout after seeing the
result.

It is to change the theory.

That negative test is therefore part of the theorem architecture:
native selection and quadratic concentration must remain distinct.

\subsection{Current theorem frontier}

The finite internal status is now:

\[
\text{quadratic contrast concentration}
\quad
\text{proved},
\]

\[
\text{finite-stage quadratic non-erasure}
\quad
\text{proved},
\]

\[
\text{native relational loading }8\to2\to1
\quad
\text{proved},
\]

\[
\text{registered-future saturation}
\quad
\text{proved},
\]

\[
\text{five-address projection-relative saturation}
\quad
\text{proved},
\]

\[
\Delta^{\mathsf T}\Delta=\frac34H
\quad
\text{proved},
\]

\[
\text{native loading contrast antimonotonicity}
\quad
\text{proved},
\]

\[
\text{global inter-face sigma}
\quad
\text{open},
\]

\[
\text{clean-release composition law}
\quad
\text{open},
\]

\[
\text{physical saturation bridge}
\quad
\text{open}.
\]

The native finite relational-saturation theorem is therefore closed.

The remaining frontier concerns stronger composition laws and physical
embedding, not the existence of the finite saturation mechanism.
